% Source-aligned English for printed p.130.
\SourcePage{130}

\LayerTarget{tohoku:section:1.5:p3}
\noindent manageable existence conditions for ``injective'' and
``projective'' objects (see No.~10). To avoid repetition, we shall first study
a very general and widely used procedure for forming new categories by means
of diagrams.

\LayerTarget{tohoku:section:1.6}
\noindent\textbf{1.6. Categories of diagrams and permanence properties.}
A \textit{diagram scheme} is a triple $(I,\Phi,d)$ consisting of two sets $I$
and $\Phi$ and a map $d$ from $\Phi$ to $I\times I$. The elements of $I$ are
the \textit{vertices}, the elements of $\Phi$ are the \textit{arrows} of the
diagram, and if $\varphi$ is an arrow of the diagram, $d(\varphi)$ is called
its \textit{direction}, characterized by the \textit{origin} and the
\textit{endpoint} of the arrow (which are therefore vertices of the scheme).

\LayerTarget{tohoku:def:composite-arrow}
A \textit{composite arrow} with origin $i$ and endpoint $j$ is, by definition,
a nonempty finite sequence of arrows of the diagram, the origin of the first
being $i$, the endpoint of each being the origin of the next, and the endpoint
of the last being $j$.

\LayerTarget{tohoku:def:diagram-category}
If $C$ is a category, a \textit{diagram in $C$ with scheme $S$} is a function
$D$ that assigns to every $i\in I$ an object $D(i)\in C$, and to every arrow
$\varphi$ with origin $i$ and endpoint $j$ a morphism $D(\varphi)$ from $D(i)$
to $D(j)$. The class of such diagrams will be denoted by $C^S$ and regarded as
a category: a morphism from $D$ to $D'$ is a family of morphisms
$v_i:D(i)\longrightarrow D'(i)$ such that, for every arrow $\varphi$ with
origin $i$ and endpoint $j$, the following diagram commutes:

\LayerTarget{tohoku:diagram:diagram-morphism-square}
\[
\begin{array}{ccc}
D(i) & \stackrel{v_i}{\longrightarrow} & D'(i) \\
{\scriptstyle D(\varphi)}\,\big\downarrow &&
  \big\downarrow\,{\scriptstyle D'(\varphi)} \\
D(j) & \stackrel{v_j}{\longrightarrow} & D'(j)
\end{array}
\qquad .
\]

Morphisms of diagrams compose in the evident way, and the axioms of a
category are verified trivially.

\LayerTarget{tohoku:def:commutative-diagram}
If $D$ is a diagram with scheme $S$, then for every composite arrow
$\varphi=(\varphi_1,\ldots,\varphi_k)$ in $S$ we define
$D(\varphi)=D(\varphi_k)\cdots D(\varphi_1)$; this is a morphism from $D(i)$
to $D(j)$ if $i$ and $j$ are respectively the origin and endpoint of
$\varphi$. The diagram $D$ is said to be \textit{commutative} if
$D(\varphi)=D(\varphi')$ whenever $\varphi$ and $\varphi'$ are two composite
arrows having the same origin and the same endpoint.

\LayerTarget{tohoku:def:commutation-relations}
More generally, let $R$ be a set consisting of pairs $(\varphi,\varphi')$ of
composite arrows having the same origins and the same endpoints, together
with composite arrows whose origin equals their endpoint. We consider the
subcategory $C^{S,R}$ of $C^S$ consisting of diagrams that satisfy the
\textit{commutation relations} $D(\varphi)=D(\varphi')$ for
$(\varphi,\varphi')\in R$, and for which $D(\varphi)$ is the identity morphism
of $D(i)$ if $\varphi\in R$ has $i$ as its origin and endpoint.

\LayerTarget{tohoku:def:additive-diagram-relations:p1}
We must also consider other types of commutation for diagrams, whose nature
varies with the categories under consideration. What follows seems to cover
the most important cases. For every $(i,j)\in I\times I$, choose a set
$R_{ij}$ of formal integral linear combinations of composite arrows with
origin $i$ and endpoint $j$, and, if $i=j$, of an auxiliary element $e_i$.
If $D$ is then a diagram with values in an additive category $C$, for every
$L\in R_{ij}$ we may define the morphism
$D(L):D(i)\longrightarrow D(j)$ by replacing in the expression for $L$ each
composite arrow $\varphi$ by $D(\varphi)$ and $e_i$ by the identity morphism
of $D(i)$. If $R$ denotes the union of the

\TohokuContinuation{131}
